\chapter{Family--defect cubic-root variational action gate}

Предыдущий gate построил flat connection с правильной three-cycle
holonomy, но оставил relation между connection и projector как constitutive
ansatz. Здесь эта relation получается как Euler--Lagrange equation одного
положительного boundary action. Цена результата формулируется заранее:
рассматривается уже выбранный nonzero condensed unit-winding sector, тогда
как происхождение самого condensate и его cubic-root bundle остаётся
отдельным parent problem.

\section{Cubic-root frame над winding sector}

Пусть charge-two pairing field не обращается в нуль на boundary circle
\(\gamma\) и задаёт фазу
\begin{equation}
 z(s)=\frac{\Phi(s)}{|\Phi(s)|}=e^{i\varphi(s)},
 \qquad
 \varphi(s+L)-\varphi(s)=2\pi\nu,
 \qquad
 \nu=\pm1.
 \label{eq:cubic-root-phase-winding}
\end{equation}
Для projector-selected field \(H\) введём уже построенный generator
\(\Omega(H)=*(u\wedge\operatorname{diag}H)\) и frame на universal cover:
\begin{equation}
 W(s)=\exp\!\left[\frac{\varphi(s)}{3}\Omega(H)\right]
 \in SO(4).
 \label{eq:cubic-root-frame}
\end{equation}
Поскольку все powers одного generator коммутируют,
\begin{equation}
 W(s)^3=\exp[\varphi(s)\Omega(H)].
\end{equation}
Frame не обязан быть single-valued как ordinary \(SO(4)\)-function. Вместо
этого он является section cubic-root bundle с transition
\begin{equation}
 W(s+L)=W(s)Z_{H,\nu},
 \qquad
 Z_{H,\nu}=\exp\!\left[\nu\frac{2\pi}{3}\Omega(H)\right],
 \qquad
 Z_{H,\nu}^{,3}=I_4.
 \label{eq:cubic-root-monodromy}
\end{equation}
Таким образом, коэффициент \(1/3\) не является fitted continuous weight:
это discrete root degree, согласованный с трёхмерным complement выбранного
rank-one projector.

Архитектура flat connections с nontrivial root holonomy имеет прямой
аналог в \(\mathbb Z_n\)-gauge constructions, где \(nA\) может быть pure
gauge при нетривиальной holonomy самого \(A\); см.
\path{arXiv:1404.3230}. Non-Abelian vortex realizations с
\(\mathbb Z_3\) Aharonov--Bohm holonomy также известны, см.
\path{arXiv:2210.11492}. Эти источники подтверждают допустимость root-bundle
архитектуры, но specific map \(H\mapsto\Omega(H)\) остаётся внутренним
результатом проекта.

\section{Единый положительный action}

Положим
\begin{equation}
 Q(H)=H^2-\frac1{\sqrt3}H-\frac14I_4,
 \qquad
 D_sW=\partial_sW+\mathcal A_sW,
\end{equation}
где \(\mathcal A_s\in\mathfrak{so}(4)\). Рассмотрим functional
\begin{equation}
 S_{\rm root}[H,W,\mathcal A]
 =\int_0^L ds\left[
  \frac1L\Tr Q(H)^2
  +L\Tr\!\left((D_sW)^T D_sW\right)
 \right].
 \label{eq:cubic-root-parent-action}
\end{equation}
Оба summands используют один и тот же representation trace; множители
\(L^{\mp1}\) только согласуют boundary dimensions. Independent multiplier
или coefficient, подбираемый по holonomy, отсутствует.

Так как \(W\in SO(4)\), matrix
\begin{equation}
 B_s=(\partial_sW)W^T
\end{equation}
лежит в \(\mathfrak{so}(4)\). Connection-dependent часть action равна
\begin{equation}
 L\Tr\!\left[(\mathcal A_s+B_s)^T(\mathcal A_s+B_s)\right].
\end{equation}
Поэтому variation по произвольному
\(\delta\mathcal A_s\in\mathfrak{so}(4)\) даёт
\begin{equation}
 \frac{\delta S_{\rm root}}{\delta\mathcal A_s}=0
 \quad\Longleftrightarrow\quad
 \mathcal A_s=-B_s=-\partial_sW\,W^T.
 \label{eq:connection-euler-lagrange}
\end{equation}
Для frame~\eqref{eq:cubic-root-frame} получаем искомую constitutive law
\begin{equation}
 \boxed{
 \mathcal A_s
 =-\frac{\partial_s\varphi}{3}\Omega(H)}.
 \label{eq:derived-projector-connection}
\end{equation}
При uniform representative
\(\partial_s\varphi=2\pi\nu/L\) это в точности relation
\begin{equation}
 \mathcal A
 =-\frac{2\pi\nu}{3L}\Omega(H)\,ds,
\end{equation}
оставленная открытой в flat-holonomy gate.

\begin{theorem}[Variational projector--connection relation]
В фиксированном nonzero unit-winding root sector canonical connection
\(\mathcal A(H,\nu)\) является единственным stationary point action
\eqref{eq:cubic-root-parent-action} по connection при фиксированном frame.
На projector branch \(Q(H)=0\) одновременно выполнено \(D_{\mathcal A}W=0\),
поэтому полный action достигает global minimum \(S_{\rm root}=0\).
\end{theorem}

Никакой дополнительной \(H\)-equation от kinetic term на этом saddle не
возникает: сам covariant derivative равен нулю. Следовательно,
projector-supercurvature equation и connection equation stationary
совместно, а identity \([\Omega(H),H]=0\) сохраняет constrained
superconnection compatibility.

\section{Holonomy, profile independence и Hessian}

Все значения \(\mathcal A_s\) пропорциональны одному generator, поэтому
path ordering исчезает и
\begin{align}
 \operatorname{Hol}_\gamma(\mathcal A)
 &=\exp\!\left[-\frac{\varphi(L)-\varphi(0)}{3}\Omega(H)\right]\\
 &=\exp\!\left[-\nu\frac{2\pi}{3}\Omega(H)\right]
 =C_{H,\nu}.
 \label{eq:root-action-holonomy}
\end{align}
Результат зависит только от winding и не зависит от локального phase
profile. Holonomy connection является inverse transition cubic-root frame,
как и требуется для parallel section \(D_{\mathcal A}W=0\).

В orthonormal basis шести generators \(\mathfrak{so}(4)\) Hessian
connection sector равен \(2I_6\). Следовательно, после фиксации frame
stationary connection является strict local minimum. Если \(W\) также
варьируется, остаётся обычная gauge orbit pure-gauge descriptions, а не
новая physical flat direction.

Машинная проверка охватывает четыре projectors, два winding sectors и два
phase profiles. Для всех шестнадцати cases получены:
\begin{align}
 \max\|D_{\mathcal A}W\|&<5.6\times10^{-16},\\
 \max\|Z_{H,\nu}^3-I\|&<5.1\times10^{-15},\\
 \max\|\operatorname{Hol}(\mathcal A)-C_{H,\nu}\|
 &<3.7\times10^{-15},\\
 \lambda_{\min}(\operatorname{Hess}_{\mathcal A}S_{\rm root})
 &=2-5\times10^{-16}.
\end{align}

\section{Точный статус и оставшийся parent gap}

Constitutive-action criterion предыдущей главы теперь проходит в
фиксированном condensed unit-winding sector: connection не назначается и
не навязывается multiplier, а следует из variation одного заранее
фиксированного sum-of-squares action с общим trace.

Однако это ещё не full boundary-superconnection closure. Action принимает
root frame \(W\) и nonzero phase \(z=\Phi/|\Phi|\) как данные выбранного
topological sector. Необходимо отдельно вывести:
\begin{enumerate}
 \item почему pairing field имеет nonzero boundary condensate;
 \item почему admissible lift является именно cubic-root bundle;
 \item как \(W\), \(H\) и sterile root line входят в один finite graded
 Hilbert module и одну superconnection;
 \item сохраняется ли механизм после BV/fluctuation audit и во втором
 normalization-sensitive sector.
\end{enumerate}

Следующий tetrahedral residual-bundle gate закрывает второй пункт на
group-theoretic уровне. Projector \(P_a\) внутри orientation-preserving
tetrahedral carrier имеет stabilizer
\(\operatorname{Stab}_{A_4}(P_a)\simeq\mathbb Z_3\), причём его два
nontrivial elements являются ровно \(C_{a,\pm}\). Поэтому cubic degree не
является новым дискретным input. Открытым остаётся gauge/global fork:
нужно вывести gauged или boundary-framed tetrahedral carrier и связать его
с nonzero ordinary charge-two condensate одним parent action.

\begin{nogocriterion}[Root-bundle embedding gate]
Нельзя объявлять family--defect parent theory построенной, пока cubic-root
frame и его condensate только постулируются как boundary sector. Следующий
pass требует вывести этот sector из representation, charge assignment,
potential и global bundle data одного parent action до обращения к family
или neutrino observables.
\end{nogocriterion}

Итак, динамическая зависимость connection от projector закрыта
вариационно. Открытым остаётся более глубокий вопрос происхождения самого
cubic-root carrier.

Машинный ledger сохранён в
\path{s2t_v4_family_defect_cubic_root_action_gate_results.json}; генератор
--- \path{s2t_v4_family_defect_cubic_root_action_gate.py}.